\documentclass[10pt,a4paper]{article}
\usepackage[utf8]{inputenc}
\usepackage[T1]{fontenc}
\usepackage[italian]{babel}
\usepackage[default]{sourcesanspro}
\usepackage{amsmath,microtype,graphicx,booktabs,tabularx,array,enumitem}
\usepackage[table]{xcolor}
\usepackage{geometry,fancyhdr,lastpage,listings,hyperref}
\geometry{left=17mm,right=17mm,top=13mm,bottom=13mm,
  headheight=18pt,headsep=4mm,footskip=8mm,includeheadfoot}
\definecolor{Navy}{HTML}{203446}
\definecolor{Muted}{HTML}{58636C}
\definecolor{Mist}{HTML}{F3F4F5}
\definecolor{Rule}{HTML}{C8CDD1}
\color{black}
\setlength{\parindent}{0pt}
\setlength{\parskip}{3pt}
\setlength{\emergencystretch}{1.5em}
\linespread{1.04}
\setlist[itemize]{leftmargin=4mm,itemsep=1pt,topsep=2pt,parsep=0pt}
\setlength{\tabcolsep}{5pt}
\renewcommand{\arraystretch}{1.12}
\hypersetup{colorlinks=true,linkcolor=Navy,urlcolor=Navy,pdflang={it-IT},
  pdftitle={Decentralised Lending Service - Relazione di progetto},
  pdfauthor={Marco Pernisco},pdfsubject={P2P Systems and Blockchains Lab 2025/26}}
\pagestyle{fancy}\fancyhf{}
\fancyhead[L]{\fontsize{8}{10}\selectfont\color{Muted}DECENTRALISED LENDING SERVICE}
\fancyhead[R]{\fontsize{8}{10}\selectfont\color{Muted}RELAZIONE DI PROGETTO}
\fancyfoot[L]{\fontsize{8}{10}\selectfont\color{Muted}Marco Pernisco / 683674}
\fancyfoot[R]{\fontsize{8}{10}\selectfont\color{Muted}\thepage\ /\ \pageref*{LastPage}}
\renewcommand{\headrulewidth}{0.3pt}
\newcommand{\entity}[1]{{\small\texttt{#1}}}
\newcommand{\pagetitle}[1]{{\fontsize{17}{20}\selectfont\bfseries\color{Navy}#1\par}\vspace{2mm}}
\newcommand{\subhead}[1]{\par\addvspace{2mm}{\fontsize{11.5}{14}\selectfont\bfseries\color{Navy}#1\par}\vspace{0.3mm}}
\lstnewenvironment{codeblock}{\lstset{
  basicstyle=\ttfamily\fontsize{9}{10.7}\selectfont,
  backgroundcolor=\color{Mist},frame=l,rulecolor=\color{Rule},framerule=0.5pt,
  xleftmargin=3mm,xrightmargin=1mm,framexleftmargin=1.5mm,framesep=1.5mm,
  aboveskip=4pt,belowskip=4pt,columns=fullflexible,keepspaces=true,
  breaklines=true,showstringspaces=false}}{}

\begin{document}
\begingroup
\centering
\setlength{\parskip}{1pt}
\includegraphics[width=36mm]{pisa_logo.png}\par
\vspace{1mm}
{\fontsize{20}{23}\selectfont\bfseries\color{Navy}Decentralised Lending Service\par}
\vspace{0.5mm}
{\small Marco Pernisco / Matricola 683674\quad
\textcolor{Muted}{P2P Systems and Blockchains Lab 2025/26}\par}
{\footnotesize\textbf{Codice sorgente:}
\url{https://github.com/Marcop-00/Decentralised-Lending-Service}\par}
\endgroup
\vspace{1mm}
\subhead{1. Architettura e protocollo}
Il servizio raccoglie ETH in un fondo comune e affida ai partecipanti il voto
sui prestiti. Il \textbf{contributore} deposita, vota, preleva capitale disponibile
e riscuote interessi o compensazioni. Il \textbf{richiedente} propone le condizioni,
risolve la proposta e, se approvata, riceve e rimborsa ETH. Un account può avere entrambi i ruoli.

\textbf{Componenti on-chain.} \entity{FundingPoolProxy} conserva ETH e stato;
\entity{FundingPool} fornisce la logica UUPS, aggiornabile solo dal proprietario.
I test verificano che lo stato sia preservato. Ogni prestito approvato crea un
\entity{LoanContract} distinto, con debitore, condizioni e rimborsi propri.
Aggiornare il pool non cambia il bytecode dei prestiti già creati.
\entity{BitcoinOracle} memorizza saldi in satoshi ed emette richieste di aggiornamento
a pagamento; il pool ne legge il saldo quando risolve una proposta.

\textbf{Componenti off-chain.} L'oracolo Python fidato elabora i primi
\textbf{131.000 blocchi Bitcoin mainnet}, uno alla volta, ricostruendo l'indice
UTXO \textbf{in memoria}, poi risponde alle richieste Ethereum. Saldi Bitcoin e
UTXO non sono salvati su disco: ogni riavvio ripete la scansione. Il voter segue
gli eventi di proposta e vota sempre YES, anche con capitale tutto bloccato. SQLite conserva
eventi e transazioni firmate, non saldi Bitcoin, per riprendere dopo errori RPC temporanei.

\textbf{Modello di fiducia.} Non si prova la proprietà dei BTC e non li si blocca.
La liquidità storica Bitcoin è valutata al cambio fisso \textbf{1 BTC = 30 ETH}.
Il parser controlla collegamento dei blocchi, hash e radici
Merkle, esclude la ricompensa non spendibile del genesis e rifiuta dati incompleti;
non è un nodo che verifica l'intero consenso Bitcoin.

\subhead{Votazione e finanziamento}
I depositi devono superare \textbf{100.000 wei}. Anche chi ha capitale bloccato
può votare, ma la risoluzione pesa il voto sul \textbf{saldo disponibile corrente}.
L'astensione vale NO; la parità respinge. Solo il richiedente originario può
risolvere, dopo almeno \textbf{12 blocchi}. Si può votare fino alla risoluzione:
\entity{votingDeadlineBlock} indica il primo blocco utile, non la fine del voto.
Tutte le durate sono differenze di altezza. La proposta specifica capitale,
tasso di interesse (1--100), durata e indirizzo BTC.

ETH disponibili o liquidità Bitcoin insufficienti chiudono e respingono la
proposta. L'approvazione chiude il voto e blocca quote proporzionali del capitale
disponibile dei contributori. I resti della divisione intera riducono il
\textbf{capitale effettivamente erogato}, senza assegnarli arbitrariamente.
La percentuale di collateral resta fissata per quel prestito alla creazione.

\subhead{Rimborso, interessi e compensazione}
Il capitale è \textbf{disponibile} o \textbf{bloccato}; interessi da riscuotere
e riserva di compensazione sono saldi separati, non prestabili. Solo il debitore
originario può rimborsare, anche dopo il default. Ogni rata copre capitale e
interessi in proporzione al debito totale; si può saldare tutto in un pagamento.

Il \textbf{rimborso del capitale} sblocca fondi nel pool, non li invia al wallet.
Si restituisce prima a chi aveva la quota \emph{iniziale} maggiore; a parità,
vince l'indirizzo numericamente maggiore. L'ordine si calcola una sola volta.
Gli \textbf{interessi} sono divisi usando il collateral fisso del prestito:
una parte entra nella riserva; il resto è un guadagno proporzionale alle quote
iniziali, riscosso con \entity{claimInterestGain()}. I calcoli cumulativi rendono
coerenti gli arrotondamenti di capitale, interessi e collateral tra le rate.
Eccedenze oltre il debito totale e resti della distribuzione degli interessi
vanno alla riserva. A pagina 4 è riportato un esempio numerico.

\textbf{Successo e default.} Il saldo completo entro il blocco di scadenza ha
successo; un prestito scaduto e non saldato è segnato fallito da una transazione,
non dal solo passare del tempo. Il fallimento si conta una volta e non è annullato
da rimborsi tardivi. Il collateral globale parte dal \textbf{50\%}, resta tra
\textbf{1 e 100\%}, scende di cinque punti per successo e sale di cinque per
fallimento; i prestiti esistenti mantengono la propria percentuale.

La \textbf{compensazione} copre solo capitale non restituito, nei limiti della
riserva. Se vuota, la richiesta registra comunque il default, paga zero e
conserva il credito per futuri reintegri. Quanto pagato riduce di pari importo
il diritto a recuperare capitale dal debitore, senza cambiare i pesi degli
interessi. Solo il proprietario può \textbf{terminare} il servizio, senza prestiti
attivi; prelievi, interessi e recuperi restano possibili.

\newpage
\pagetitle{2. Sicurezza e limiti operativi}
\subhead{Reentrancy: modifica e attacco eseguibile}
\entity{FundingPoolVulnerable}, usato solo nei test, eredita dal pool reale e
sostituisce il prelievo con scritture di saldi memorizzati \textbf{dopo} il
trasferimento esterno:
\begin{codeblock}
uint256 beforePosition = totalContributed[msg.sender];
uint256 beforePool = totalUnlockedFunds;
(bool sent, ) = payable(msg.sender).call{value: amount}("");
if (!sent) revert TransferFailed();
totalContributed[msg.sender] = beforePosition - amount;
totalUnlockedFunds = beforePool - amount;
\end{codeblock}
\entity{ReentrancyAttacker} richiama ricorsivamente \entity{withdraw(amount)}
dalla funzione di ricezione. In \entity{test/reentrancy.test.ts}, un deposito
di 1 ETH dell'attaccante permette di sottrarre i 10 ETH del contributore onesto
dal proxy vulnerabile. Lo stesso attaccante non ruba quei 10 ETH dal proxy
normale, che aggiorna la contabilità prima del trasferimento. Una semplice
sottrazione dal saldo corrente dopo la chiamata potrebbe invece causare un
revert per underflow in Solidity 0.8; le assegnazioni di valori ormai obsoleti
rendono l'esempio davvero sfruttabile. Le varianti di test non sono mai
distribuite dallo setup del servizio. Queste prove non garantiscono la sicurezza generale.

\subhead{Guadagno scorretto senza reentrancy}
Un contributore di maggioranza con \textbf{6 ETH} può approvare un prestito da
\textbf{10 ETH} a un debitore complice, finanziato anche dai \textbf{4 ETH}
di un partecipante onesto, e poi non rimborsarlo. Il gruppo colluso guadagna
\textbf{4 ETH netti}, non 10, prima di gas e altri flussi. Può indicare qualsiasi
indirizzo BTC abbastanza ricco, perché non si prova la proprietà. Anche la
priorità dei rimborsi favorisce le quote grandi: con quote 6/4, un rimborso di
capitale pari a 6 ripristina la posizione maggiore prima che il contributore
minore recuperi capitale. Sono conseguenze delle regole economiche del protocollo,
non di una vulnerabilità reentrancy.

\textbf{Limiti.} Le posizioni vuote vengono rimosse, ma il lavoro cresce con il
numero di contributori e deve rientrare nel limite di gas del blocco. I checkpoint
rilevano la sostituzione della chain; i conflitti di nonce restano ispezionabili.
Il recupero dalle riorganizzazioni non è automatico.

\newpage
\pagetitle{3. Esecuzione del servizio}
Eseguire i comandi dalla \textbf{radice del progetto}, salvo diversa indicazione.
Servono una versione Node pari supportata (22.13.0+), Python 3.10+ e
\textbf{Geth 1.13.14}, client Clique precedente a PoS compatibile con il genesis della rete privata.

\subhead{1. Installare le dipendenze e validare i dati Bitcoin}
\begin{codeblock}
./setup.sh
./.venv/bin/python python/oracle/check_data.py --data-dir bitcoin --summary
\end{codeblock}
Lo setup installa le dipendenze Node bloccate e i pacchetti Python, crea
\entity{.venv}, compila i contratti e scarica Geth. Installazione e download
del compilatore non già in cache richiedono internet; saltare lo setup se già
completato. Fornire file mainnet \entity{blk*.dat} decodificati dal genesis
all'altezza 130.999. Il controllo offline non richiede chiavi o Geth: il successo
riporta \entity{valid: true}, \textbf{131.000 blocchi}, altezza \textbf{130.999}
e indirizzi con fondi. Non sostituire dati mancanti o incompleti con saldi fittizi.

\subhead{2. Avviare la blockchain privata / terminale 1}
\begin{codeblock}
cd private-chain
./geth --datadir=./data init project2526genesis.json
./geth --datadir=./data --keystore=./keystore --networkid=202526 \
  --http --http.addr=127.0.0.1 --http.port=8545 --http.api=eth,net,web3 \
  --mine --miner.gasprice=100000000 \
  --miner.etherbase=d278d247a52c550508ea2b2c9321d816238fb523 \
  --unlock d278d247a52c550508ea2b2c9321d816238fb523 \
  --password password.txt --allow-insecure-unlock --ipcdisable
\end{codeblock}
Lasciare il nodo attivo.
\subhead{3. Creare gli account e distribuire / terminale 2, radice del progetto}
\begin{codeblock}
./.venv/bin/python python/setup.py
\end{codeblock}
L'account prefunded fornito \textbf{trasferisce solo ETH} a tre account appena
generati. Il nuovo deployer crea oracolo, implementazione e proxy inizializzato.
Lo setup ricompila il bytecode ordinario, verifica chain ID, ricevute, codice e
inizializzazione, poi scrive configurazioni protette con \entity{DEPLOYMENT\_BLOCK}.
Ripeterlo crea una distribuzione separata.

\subhead{4. Avviare oracolo e voter / terminali separati}
\begin{codeblock}
./.venv/bin/python python/oracle/oracle_service.py
./.venv/bin/python python/voter/auto_voter.py
\end{codeblock}
Attendere \textbf{ORACLE READY}: l'oracolo deve terminare la scansione Bitcoin
iniziale in memoria. Non c'è un'attesa fissa da configurare; ogni riavvio ripete
la scansione. Lasciare attivi entrambi i servizi durante la demo e conservare
le code SQLite dopo errori RPC temporanei: sono separate dall'indice Bitcoin.

\subhead{5. Simulare rimborso e default / altro terminale, radice del progetto}
\begin{codeblock}
./.venv/bin/python python/simulate.py \
  --btc-address 1A1zP1eP5QGefi2DMPTfTL5SLmv7DivfNa --principal 1
\end{codeblock}
Questo indirizzo ha \textbf{4.000.000 satoshi} nello snapshot validato, sufficienti
per 1 ETH al cambio fisso. Lo scenario A
rimborsa 1 ETH al 5\% in due rate e riscuote il guadagno; B lascia insoluti 0,5 ETH,
richiede compensazione e riprova a riserva vuota. Entrambe le proposte attendono
voti YES e 12 blocchi reali. Saldi, stati e controlli contabili sono stampati dopo
le modifiche significative.

\textbf{Tempi e ripresa.} I blocchi Clique richiedono circa 10 secondi: prevedere
alcuni minuti per le due votazioni. \entity{--timeout} imposta i secondi massimi
\emph{per ogni fase asincrona} (300 predefiniti). In caso di timeout, controllare
nodo e log dei servizi prima di riprovare; non eliminare chain, chiavi,
configurazioni o code. La contabilità è spiegata a pagina 4; riusare una
distribuzione può produrre numeri diversi.

\newpage
\pagetitle{4. Simulazione: rimborso e default}
L'esecuzione locale osservata usa \entity{--principal 1}. L'oracolo risponde con
\textbf{4.000.000 satoshi = 0,04 BTC}, pari a \textbf{1,2 ETH} di capacità al
cambio di 30 ETH/BTC. I BTC non sono trasferiti né bloccati. Il contributore
deposita \textbf{50 ETH}: il wallet passa da 100 a circa 50 ETH dopo il gas,
mentre il pool contiene 50 ETH di capitale disponibile.

\subhead{Primo prestito: 1 ETH, interesse al 5\%, due rate}
Dopo il voto YES e 12 blocchi, l'approvazione crea e finanzia un contratto
dedicato. Il pool contiene 49 ETH; il contributore ha \textbf{49 disponibili
+ 1 bloccato}. Il debito è \textbf{1,05 ETH}. Con \textbf{collateral al 50\%},
ogni rata da 0,525 ETH restituisce 0,5 ETH di capitale e paga 0,025 ETH di
interessi, divisi a metà tra riserva e guadagno del contributore:

\begin{tabularx}{\linewidth}{@{}Xrr@{}}
\toprule
\textbf{Destinazione} & \textbf{Ogni rata (ETH)} & \textbf{Totale (ETH)}\\\midrule
Capitale ripristinato come disponibile & 0,5000 & 1,0000\\
Interessi accreditati alla riserva & 0,0125 & 0,0250\\
Interessi accreditati come guadagno & 0,0125 & 0,0250\\\midrule
Pagamento del debitore & 0,5250 & 1,0500\\\bottomrule
\end{tabularx}

Il prestito diventa \entity{Repaid}: il capitale disponibile torna a 50 ETH e
quello bloccato a zero. Il contributore riscuote 0,025 ETH di guadagno; nel pool
restano \textbf{50 ETH di capitale + 0,025 ETH di riserva}. Il collateral globale
scende dal 50 al \textbf{45\%}; quello del primo prestito resta al 50\%.

\subhead{Secondo prestito: la scadenza non esegue transazioni}
Una nuova proposta riceve YES ed è risolta dopo altri 12 blocchi. Eroga
\textbf{0,5 ETH} al 10\%, con collateral fisso al \textbf{45\%}, e non viene
rimborsata intenzionalmente. Il finanziamento avviene al blocco \textbf{328},
la scadenza è \textbf{331}. Al 332 il prestito è scaduto ma resta memorizzato
come \entity{Funded}: il solo avanzare dell'altezza non esegue codice.
La transazione di compensazione al blocco \textbf{333} registra
\entity{Defaulted}, azzera i prestiti attivi, conta un fallimento e riporta
il collateral globale al \textbf{50\%}. Quello del secondo prestito resta al 45\%.

\subhead{Compensazione: riserva limitata e credito residuo}
Il capitale insoluto è 0,5 ETH, ma la riserva contiene solo 0,025 ETH.
Il contributore riceve \textbf{0,025 ETH direttamente nel wallet}; il credito
residuo è \textbf{0,475 ETH}. Il tentativo al blocco \textbf{334} paga zero
senza cancellare il credito né contare il fallimento una seconda volta.

\begin{tabularx}{\linewidth}{@{}Xrr@{}}
\toprule
\textbf{Contabilità del pool (ETH)} & \textbf{Prima} & \textbf{Dopo / nuovo tentativo}\\\midrule
Posizione totale del contributore & 50,000 & 49,975\\
Capitale disponibile del contributore & 49,500 & 49,500\\
Capitale bloccato del contributore & 0,500 & 0,475\\
Riserva di compensazione & 0,025 & 0\\
ETH effettivamente detenuti dal pool & 49,525 & 49,500\\\bottomrule
\end{tabularx}

\textbf{Interpretazione.} I 0,475 ETH bloccati sono un \textbf{credito non
recuperato, non denaro prelevabile}. La compensazione riduce sia la posizione
totale sia il capitale bloccato, lasciando invariato il disponibile:
$49{,}975-0{,}475=49{,}5$ ETH. Reintegri della riserva o rimborsi tardivi possono
consentire altri recuperi, ma non sono garantiti. \textbf{Zero prestiti attivi
non significa zero perdite.}

\subhead{Saldi dei wallet e numero di transazioni}
Il wallet finale del contributore contiene \textbf{50,0499454037 ETH}, con il gas
già considerato: vi sono arrivati guadagno e compensazione, mentre 49,5 ETH
restano disponibili nel pool. Questa esecuzione aggiunge \textbf{11 transazioni
della demo + 1 aggiornamento dell'oracolo + 2 voti YES = 14}; lo setup ne aggiunge
sei (tre trasferimenti e tre distribuzioni). L'oracolo contiene 0,0000098852 ETH
perché anche una precedente richiesta per un indirizzo inadeguato aveva pagato
la commissione; quella coppia richiesta/risposta non rientra nelle 14 transazioni.
Blocchi e saldi sono quelli dell'esecuzione osservata; riusare una distribuzione
o cambiare dati può modificarli.

\newpage
\pagetitle{5. Test e valutazione del gas}
\subhead{5.1 Esecuzione e risultati dei test}
Dopo lo setup delle dipendenze (pagina 3), lavorare dalla \textbf{radice del progetto}.
\textbf{Non servono Geth, oracolo, voter, chiavi, configurazioni dei servizi o
file Bitcoin reali.} Hardhat usa chain simulate isolate; Python usa dati
temporanei e servizi simulati con mock. I test non modificano la distribuzione
privata esistente. Non eseguire \entity{python/setup.py} soltanto per fare i test.

\subhead{Compilazione e suite complete}
\begin{codeblock}
npm run typecheck
npm test
./.venv/bin/python -m unittest discover -s python/tests -v
\end{codeblock}
\textbf{Typecheck:} compila i contratti ed esegue \entity{tsc --noEmit}; basta
l'uscita senza errori, anche se TypeScript non stampa messaggi.
\textbf{Contratti:} l'esito atteso è \textbf{34 passing}, cioè \textbf{4 test
Solidity + 30 Node.js/TypeScript}. Un test Solidity usa \textbf{256 input
generati} mediante fuzzing: sono ripetizioni dello stesso test, non 256 test
aggiuntivi. \textbf{Python:} attendersi \textbf{Ran 49 tests} e \textbf{OK}.

Gli scenari dei contratti coprono vincoli del voto e variazioni di saldo, soglie
BTC, resti di allocazione, rate, ordine dei rimborsi, interessi, default,
compensazioni ripetute, autorizzazioni, terminazione, aggiornamenti e reentrancy.
Controlli espliciti verificano emittenti e parametri degli eventi e contratti
indipendenti per due debitori. I test Python coprono parser, esclusione del
genesis, UTXO, indicizzazione atomica, ripresa di eventi e transazioni, voter
sempre favorevole, protezioni dello setup e compensazione della demo. I saldi
simulati verificano la logica, non validano il dataset Bitcoin reale.

\subhead{Eseguire un controllo mirato}
\begin{codeblock}
npm test -- test/advanced_behaviors.test.ts
npm test -- test/events_and_isolation.test.ts
npm test -- test/reentrancy.test.ts
npm test -- test/upgrade.test.ts
npm test -- contracts/test/LoanContract.t.sol
./.venv/bin/python -m unittest discover -s python/tests \
  -p 'test_demo_scenarios.py' -v
\end{codeblock}
Scegliere il comando per: contabilità; eventi e indipendenza dei prestiti;
reentrancy; aggiornamenti; regole native del prestito; oppure \textbf{sei
controlli offline della demo}. Gli eventi includono un nuovo tentativo di
compensazione che non deve emettere un secondo default. Il test di indipendenza
rimborsa due prestiti con debitori diversi e verifica che nessuno alteri
condizioni, rimborsi o quote dell'altro.

\subhead{Esaminare la copertura del codice}
\begin{codeblock}
npm run test:coverage
\end{codeblock}
\textbf{Coverage:} attendersi \textbf{34 passing} e percentuali di righe e
istruzioni coperte. Aprire \entity{coverage/html/index.html} nel browser.
È copertura dei contratti, non Python né prova di sicurezza. Il bytecode
strumentato altera il gas: questa modalità lascia intenzionalmente invariate
le misure ordinarie. Comandi di profilazione e tabella sono a pagina 6;
non usare il gas della coverage per valutare i costi.

\subhead{Riconoscere il successo e risolvere i problemi}
Un test può causare un revert e passare comunque: verifica il rifiuto di
un'operazione non valida. Il test reentrancy passa solo se la variante insicura
è svuotata \emph{e} quella normale protegge i fondi onesti. Leggere il riepilogo
finale: asserzioni fallite, \entity{FAIL}/\entity{ERROR} o uscita non nulla richiedono
indagini. I conteggi descrivono questa versione e possono crescere con nuovi test.

\textbf{Moduli o interprete mancanti:} usare \entity{./.venv/bin/python} e verificare
cartella, versione Node e completamento di \entity{./setup.sh}.
\textbf{Test fallito:} rieseguire il file e leggere il primo scarto tra atteso
e ottenuto, senza eliminare dati della chain o code.
\textbf{Dati Bitcoin mancanti:} correggere il dataset per controllo e demo
(pagina 3); i mock non lo validano. La demo reale è una prova di integrazione:
non sostituisce i test automatici, né viceversa.

\newpage
\pagetitle{5. Test e valutazione del gas (continua)}
\subhead{5.2 Misurazione e interpretazione del gas}
I test funzionali verificano il comportamento; la profilazione misura il costo
delle operazioni. Entrambi usano lo stesso progetto Hardhat e gli stessi
contratti, con bytecode ordinario per le misurazioni.

\textbf{Ambiente.} Hardhat 3.18.1, TypeScript 6.0.3, \entity{node:test}, viem
2.57.3; Solidity 0.8.24, ottimizzatore con \entity{runs=200}, EVM/hardfork
\textbf{Berlin}, chain ID simulato 202526. Lo scenario ha un solo contributore.

\subhead{Riprodurre le misure}
\begin{codeblock}
npm run test:gas
\end{codeblock}
L'esito atteso è \textbf{1 passing}: un test di profilazione per \textbf{20 scenari}.
Registra \entity{gasUsed} delle ricevute, incluso il gas intrinseco, in
\entity{docs/gas-measurements.json}, creato automaticamente durante il test.
La tabella seguente riporta i risultati misurati per questa implementazione.
Anche \entity{npm test} include la profilazione e riscrive il JSON.

\subhead{Costi del gas misurati}
\begingroup\fontsize{9}{11}\selectfont
\renewcommand{\arraystretch}{1.08}
\par\noindent\begin{tabularx}{\linewidth}{@{}lXr@{}}
\toprule
\textbf{Operazione} & \textbf{Scenario} & \textbf{Unità di gas} \\
\midrule
\addlinespace[2pt]
\multicolumn{3}{l}{\cellcolor{Mist}\textbf{\color{Navy}FundingPool / tramite proxy}} \\
\texttt{contribute()} & Nuovo contributore & 140.923 \\
\texttt{withdraw()} & Prelievo parziale & 47.933 \\
\texttt{withdraw()} & Prelievo finale & 32.412 \\
\texttt{createLoanProposal()} & Nuova proposta & 234.798 \\
\texttt{voteOnProposal()} & Voto favorevole & 106.936 \\
\texttt{executeApprovedLoan()} & Proposta approvata & 1.049.274 \\
\texttt{executeApprovedLoan()} & Proposta respinta & 90.860 \\
\texttt{claimInterestGain()} & Riscossione del guadagno & 24.691 \\
\texttt{requestCompensation()} & Compensazione pagata & 52.652 \\
\texttt{requestCompensation()} & Riserva vuota & 43.786 \\
\texttt{upgradeToAndCall()} & V2 compatibile & 37.769 \\
\texttt{terminate()} & Nessun prestito attivo & 36.295 \\
\addlinespace[2pt]
\multicolumn{3}{l}{\cellcolor{Mist}\textbf{\color{Navy}LoanContract}} \\
\texttt{repay()} & Rata parziale & 168.082 \\
\texttt{repay()} & Rata finale & 90.716 \\
\texttt{repay()} & Rimborso in ritardo & 138.782 \\
\texttt{markDefaulted()} & Prestito scaduto & 103.758 \\
\addlinespace[2pt]
\multicolumn{3}{l}{\cellcolor{Mist}\textbf{\color{Navy}BitcoinOracle}} \\
\texttt{requestUpdate()} & Richiesta a pagamento & 25.193 \\
\texttt{update()} & Primo saldo non nullo & 49.426 \\
\texttt{update()} & Saldo già presente & 32.326 \\
\texttt{withdrawFees()} & Riscossione delle commissioni & 30.868 \\
\bottomrule
\end{tabularx}\par
\endgroup


\subhead{Leggere i valori e calibrare la commissione dell'oracolo}
I valori sono \textbf{unità di gas}, non wei né il limite di gas inviato.
La commissione in wei è \textbf{gas consumato per prezzo del gas}.
Le letture tramite \entity{eth\_call} non richiedono una commissione di transazione.

\textbf{Commissione dell'oracolo:}
$49\,426\times0{,}1\text{ gwei}=4\,942\,600\,000\,000$ wei.
La calibrazione usa il primo saldo non nullo di un indirizzo storico di 34
caratteri; aggiornare una voce esistente costa meno. Cambiare bytecode o
impostazioni richiede una nuova calibrazione.

\textbf{Validità delle misure.} Sono risultati riproducibili di scenari specifici,
non limiti massimi universali. Dipendono dallo stato dello storage e dal numero
di contributori: il lavoro del pool può crescere con quest'ultimo. La coverage
esegue gli stessi percorsi con bytecode strumentato senza sovrascrivere il JSON
ordinario; per valutare i costi usare il comando di profilazione riportato sopra.
\end{document}
